%% GENERATED by python/paper/build.py from results/ -- do not edit by hand.
%% Rebuild:  .venv\Scripts\python.exe python\paper\build.py --only US_OtherShocks
%% Source:   results/shocks/US_shocksCommonX.csv
\begin{table}[!htb]
\centering
\begin{threeparttable}
\caption{French income distribution and voting patterns in US}
\label{table:US:otherShocks}
\renewcommand{\arraystretch}{1.25}
\begin{tabularx}{.9\textwidth}{lYYY}
\toprule
\textbf{Scenario} & \textbf{Tax rate} & \textbf{Savings rate} & \textbf{Avg. workweek} \\
\midrule 
Baseline & 14.43\% & 15.37\% & 39.39 \\
Income distribution & 13.21\% & $+0.44$ p.p. & 40.36 \\
Voting & 15.32\% & $-0.32$ p.p. & 39.19 \\
All French characteristics & 13.85\% & $+0.21$ p.p. & 35.29 \\[.5em]\hline\\[-.75em]
France (own calibration) & 21.29\% & $-1.94$ p.p. & 35.44 \\
\bottomrule
\end{tabularx}
\begin{tablenotes}
\footnotesize
\item \textit{Note:} $\rho = 1.0$, full effect. Each row is a separate equilibrium path: the borrowed characteristics hold throughout and the economy starts from its own steady state, so the row describes a country that has always had this mix rather than the US hit by a surprise in 2020. Income distribution replaces $\eta_i$ with France\textquotesingle s while holding $X_i$ \emph{and} holding $\theta$ at the US design, so it is a change in inequality alone; pension design is the separate counterfactual of Table~\ref{table:US:pensChars}. All French characteristics imposes France\textquotesingle s $\eta_i$, voting weights $\mu_i$ and level of $X_i$ at once; the last, a pure rescaling of the hours unit, moves hours and nothing else. The last row is France\textquotesingle s own calibrated path, its savings rate reported as the distance from the US baseline. Common-$X$ calibration: see the note to Table~\ref{table:US:Calib}.
\end{tablenotes}
\end{threeparttable}
\end{table}
